% ============================================================================
% SEC03.TEX - Section 11.3: Script Pengejar (Chasing)
% ============================================================================

\section{Script Pengejar (Chasing)}

Buat script \texttt{EnemyAI.cs}:

\begin{lstlisting}[language={[Sharp]C}]
using UnityEngine;

public class EnemyAI : MonoBehaviour
{
    public float speed = 3f;
    private Transform player;

    void Start()
    {
        // Cari objek Player secara otomatis saat spawn
        GameObject playerObj = GameObject.FindGameObjectWithTag("Player");
        if (playerObj != null)
        {
            player = playerObj.transform;
        }
    }

    void Update()
    {
        if (player != null)
        {
            // Tentukan arah ke pemain
            // Arah = (Posisi Tujuan - Posisi Sekarang).normalized
            Vector3 direction = (player.position - transform.position).normalized;
            
            // Gerakkan musuh
            transform.position += direction * speed * Time.deltaTime;
            
            // Putar musuh menghadap pemain (Opsional, gunakan Atan2)
            float angle = Mathf.Atan2(direction.y, direction.x) * Mathf.Rad2Deg;
            transform.rotation = Quaternion.Euler(new Vector3(0, 0, angle - 90));
        }
    }
}
\end{lstlisting}

\begin{tips}
Jangan lupa memberi Tag \textbf{"Player"} pada object Player utama Anda! Fungsi \texttt{GameObject.\allowbreak FindGameObjectWithTag} akan gagal jika tidak ada tag tersebut.
\end{tips}
